\section{Preguntas y diseño experimental}
\label{sec:preguntas}

El artículo organiza su indagación en cinco preguntas que cubren la transferencia de dominio, la transferencia de tarea, el valor del preentrenamiento especializado, el efecto del tamaño de los datos y el cambio en las representaciones internas \cite{piperno2025}.

\begin{itemize}
  \item \textbf{Pregunta 1.} ¿Puede transferirse el conocimiento de dominio de un modelo preentrenado en ese dominio cuando no se dispone de datos etiquetados en inglés para la tarea?
  \item \textbf{Pregunta 2.} ¿Puede transferirse el conocimiento de tarea desde un modelo ajustado en inglés cuando ese tipo de datos sí existe?
  \item \textbf{Pregunta 3.} Si no hay un modelo adaptado al dominio disponible en inglés, ¿qué beneficios tiene destilar el conocimiento de tarea desde un modelo ajustado a partir del mismo punto de partida que el estudiante?
  \item \textbf{Pregunta 4.} ¿Cómo influye la cantidad de datos usada en el ajuste sobre el desempeño final?
  \item \textbf{Pregunta 5.} ¿Qué diferencias aparecen en las representaciones internas del profesor y del estudiante después de la destilación?
\end{itemize}

Las tres primeras se responden contrastando un profesor con otro. La pregunta 4 es de escala y se responde con la curva de desempeño. La pregunta 5 es de mecanismo y se responde con análisis de representaciones.

El diseño se detalla en la sección~\ref{sec:metodo}: el estudiante es mBERT, los profesores son $B_p$, $B_{ft}$ y $mB_{ft}$, y el baseline es el mismo mBERT ajustado sin destilación en la misma fracción de datos.

La condición que el artículo describe como ausencia de datos en inglés merece una aclaración. En la pregunta 1 significa que no hay etiquetas inglesas de la tarea en el ajuste final, no que el inglés desaparezca. El profesor se entrenó con texto inglés, la destilación usa oraciones inglesas y sus traducciones, y los autores reconocen que todos los modelos se entrenaron parcialmente en inglés. El método no sirve para una institución que carece por completo de recursos en inglés; sirve para una que tiene acceso a un modelo preentrenado y a un traductor, pero no tiene anotación local suficiente. La condición realista es la escasez de etiquetas, no la ausencia de inglés.
